    
    

    
    \hypertarget{a-minimal-alf-run}{%
\section{A minimal ALF run}\label{a-minimal-alf-run}}

    In this bare-bones example we use the
\href{https://git.physik.uni-wuerzburg.de/ALF/pyALF}{pyALF} interface to
run the canonical Hubbard model on a default configuration: a
\(6\times6\) square grid, with interaction strength \(U=4\) and inverse
temperature \(\beta = 5\).

Bellow we go through the steps for performing the simulation and
outputting observables.

\begin{center}\rule{0.5\linewidth}{\linethickness}\end{center}

\textbf{1.} Import \texttt{Simulation} class from the \texttt{py\_alf}
python module, which provides the interface with ALF:

    \begin{tcolorbox}[breakable, size=fbox, boxrule=1pt, pad at break*=1mm,colback=cellbackground, colframe=cellborder]
\prompt{In}{incolor}{1}{\boxspacing}
\begin{Verbatim}[commandchars=\\\{\}]
\PY{k+kn}{from} \PY{n+nn}{py\PYZus{}alf} \PY{k}{import} \PY{n}{Simulation}            \PY{c+c1}{\PYZsh{} Interface with ALF}
\end{Verbatim}
\end{tcolorbox}

    \textbf{2.} Create an instance of \texttt{Simulation}, setting
parameters as desired:

    \begin{tcolorbox}[breakable, size=fbox, boxrule=1pt, pad at break*=1mm,colback=cellbackground, colframe=cellborder]
\prompt{In}{incolor}{2}{\boxspacing}
\begin{Verbatim}[commandchars=\\\{\}]
\PY{n}{sim} \PY{o}{=} \PY{n}{Simulation}\PY{p}{(}
    \PY{l+s+s2}{\PYZdq{}}\PY{l+s+s2}{Hubbard}\PY{l+s+s2}{\PYZdq{}}\PY{p}{,}                           \PY{c+c1}{\PYZsh{} Hamiltonian}
    \PY{p}{\PYZob{}}                                    \PY{c+c1}{\PYZsh{} Model and simulation parameters for each Simulation instance}
    \PY{l+s+s2}{\PYZdq{}}\PY{l+s+s2}{Model}\PY{l+s+s2}{\PYZdq{}}\PY{p}{:} \PY{l+s+s2}{\PYZdq{}}\PY{l+s+s2}{Hubbard}\PY{l+s+s2}{\PYZdq{}}\PY{p}{,}                  \PY{c+c1}{\PYZsh{} Base model}
    \PY{l+s+s2}{\PYZdq{}}\PY{l+s+s2}{Lattice\PYZus{}type}\PY{l+s+s2}{\PYZdq{}}\PY{p}{:} \PY{l+s+s2}{\PYZdq{}}\PY{l+s+s2}{Square}\PY{l+s+s2}{\PYZdq{}}\PY{p}{\PYZcb{}}\PY{p}{,}           \PY{c+c1}{\PYZsh{} Lattice type}
\PY{p}{)}
\end{Verbatim}
\end{tcolorbox}

    \textbf{3.} Compile ALF, downloading it first from the
\href{https://git.physik.uni-wuerzburg.de:ALF}{ALF repository} if not
found locally. This may take a few minutes:

    \begin{tcolorbox}[breakable, size=fbox, boxrule=1pt, pad at break*=1mm,colback=cellbackground, colframe=cellborder]
\prompt{In}{incolor}{3}{\boxspacing}
\begin{Verbatim}[commandchars=\\\{\}]
\PY{n}{sim}\PY{o}{.}\PY{n}{compile}\PY{p}{(}\PY{p}{)}                            \PY{c+c1}{\PYZsh{} Compilation needs to be performed only once}
\end{Verbatim}
\end{tcolorbox}

    \begin{Verbatim}[commandchars=\\\{\}]
Repository /home/stafusa/ALF/pyALF/Notebooks/ALF does not exist, cloning from
git@git.physik.uni-wuerzburg.de:ALF/ALF.git
Compiling ALF{\ldots} Done.
    \end{Verbatim}

    \textbf{4.} Perform the simulation as specified in \texttt{sim}:

    \begin{tcolorbox}[breakable, size=fbox, boxrule=1pt, pad at break*=1mm,colback=cellbackground, colframe=cellborder]
\prompt{In}{incolor}{4}{\boxspacing}
\begin{Verbatim}[commandchars=\\\{\}]
\PY{n}{sim}\PY{o}{.}\PY{n}{run}\PY{p}{(}\PY{p}{)}                                \PY{c+c1}{\PYZsh{} Perform the actual simulation in ALF}
\end{Verbatim}
\end{tcolorbox}

    \begin{Verbatim}[commandchars=\\\{\}]
Prepare directory "/home/stafusa/ALF/pyALF/Notebooks/Hubbard\_Square" for Monte
Carlo run.
Create new directory.
Run /home/stafusa/ALF/pyALF/Notebooks/ALF/Prog/Hubbard.out
    \end{Verbatim}

    \textbf{5.} Perform some simple analyses:

    \begin{tcolorbox}[breakable, size=fbox, boxrule=1pt, pad at break*=1mm,colback=cellbackground, colframe=cellborder]
\prompt{In}{incolor}{5}{\boxspacing}
\begin{Verbatim}[commandchars=\\\{\}]
\PY{n}{sim}\PY{o}{.}\PY{n}{analysis}\PY{p}{(}\PY{p}{)}                           \PY{c+c1}{\PYZsh{} Perform default analysis; list observables}
\end{Verbatim}
\end{tcolorbox}

    \begin{Verbatim}[commandchars=\\\{\}]
Analysing Ener\_scal
Analysing Part\_scal
Analysing Pot\_scal
Analysing Kin\_scal
Analysing Den\_eq
Analysing SpinZ\_eq
Analysing Green\_eq
Analysing SpinXY\_eq
Analysing SpinT\_eq
Analysing SpinXY\_tau
Analysing SpinZ\_tau
Analysing Den\_tau
Analysing Green\_tau
Analysing SpinT\_tau
    \end{Verbatim}

    \textbf{6.} Store computed observables list:

    \begin{tcolorbox}[breakable, size=fbox, boxrule=1pt, pad at break*=1mm,colback=cellbackground, colframe=cellborder]
\prompt{In}{incolor}{6}{\boxspacing}
\begin{Verbatim}[commandchars=\\\{\}]
\PY{n}{obs} \PY{o}{=} \PY{n}{sim}\PY{o}{.}\PY{n}{get\PYZus{}obs}\PY{p}{(}\PY{p}{)}                      \PY{c+c1}{\PYZsh{} Dictionary for the observables}
\end{Verbatim}
\end{tcolorbox}

    which are available for further analyses. For instance, the internal
energy of the system (and its error) is accessed by:

    \begin{tcolorbox}[breakable, size=fbox, boxrule=1pt, pad at break*=1mm,colback=cellbackground, colframe=cellborder]
\prompt{In}{incolor}{7}{\boxspacing}
\begin{Verbatim}[commandchars=\\\{\}]
\PY{n}{obs}\PY{p}{[}\PY{l+s+s1}{\PYZsq{}}\PY{l+s+s1}{Ener\PYZus{}scalJ}\PY{l+s+s1}{\PYZsq{}}\PY{p}{]}\PY{p}{[}\PY{l+s+s1}{\PYZsq{}}\PY{l+s+s1}{obs}\PY{l+s+s1}{\PYZsq{}}\PY{p}{]}
\end{Verbatim}
\end{tcolorbox}

            \begin{tcolorbox}[breakable, size=fbox, boxrule=.5pt, pad at break*=1mm, opacityfill=0]
\prompt{Out}{outcolor}{7}{\boxspacing}
\begin{Verbatim}[commandchars=\\\{\}]
array([[-29.983503,   0.232685]])
\end{Verbatim}
\end{tcolorbox}
        
    \begin{center}\rule{0.5\linewidth}{\linethickness}\end{center}

\textbf{7.} Running again: The simulation can be resumed to increase the
precision of the results.

    \begin{tcolorbox}[breakable, size=fbox, boxrule=1pt, pad at break*=1mm,colback=cellbackground, colframe=cellborder]
\prompt{In}{incolor}{8}{\boxspacing}
\begin{Verbatim}[commandchars=\\\{\}]
\PY{n}{sim}\PY{o}{.}\PY{n}{run}\PY{p}{(}\PY{p}{)}
\PY{n}{sim}\PY{o}{.}\PY{n}{analysis}\PY{p}{(}\PY{p}{)}
\PY{n}{obs2} \PY{o}{=} \PY{n}{sim}\PY{o}{.}\PY{n}{get\PYZus{}obs}\PY{p}{(}\PY{p}{)}
\PY{n+nb}{print}\PY{p}{(}\PY{n}{obs2}\PY{p}{[}\PY{l+s+s1}{\PYZsq{}}\PY{l+s+s1}{Ener\PYZus{}scalJ}\PY{l+s+s1}{\PYZsq{}}\PY{p}{]}\PY{p}{[}\PY{l+s+s1}{\PYZsq{}}\PY{l+s+s1}{obs}\PY{l+s+s1}{\PYZsq{}}\PY{p}{]}\PY{p}{)}
\PY{n+nb}{print}\PY{p}{(}\PY{l+s+s2}{\PYZdq{}}\PY{l+s+se}{\PYZbs{}n}\PY{l+s+s2}{Running again reduced the error from }\PY{l+s+s2}{\PYZdq{}}\PY{p}{,} \PY{n}{obs}\PY{p}{[}\PY{l+s+s1}{\PYZsq{}}\PY{l+s+s1}{Ener\PYZus{}scalJ}\PY{l+s+s1}{\PYZsq{}}\PY{p}{]}\PY{p}{[}\PY{l+s+s1}{\PYZsq{}}\PY{l+s+s1}{obs}\PY{l+s+s1}{\PYZsq{}}\PY{p}{]}\PY{p}{[}\PY{l+m+mi}{0}\PY{p}{]}\PY{p}{[}\PY{l+m+mi}{1}\PY{p}{]}\PY{p}{,}\PY{l+s+s2}{\PYZdq{}}\PY{l+s+s2}{ to }\PY{l+s+s2}{\PYZdq{}}\PY{p}{,} \PY{n}{obs2}\PY{p}{[}\PY{l+s+s1}{\PYZsq{}}\PY{l+s+s1}{Ener\PYZus{}scalJ}\PY{l+s+s1}{\PYZsq{}}\PY{p}{]}\PY{p}{[}\PY{l+s+s1}{\PYZsq{}}\PY{l+s+s1}{obs}\PY{l+s+s1}{\PYZsq{}}\PY{p}{]}\PY{p}{[}\PY{l+m+mi}{0}\PY{p}{]}\PY{p}{[}\PY{l+m+mi}{1}\PY{p}{]}\PY{p}{,} \PY{l+s+s2}{\PYZdq{}}\PY{l+s+s2}{.}\PY{l+s+s2}{\PYZdq{}}\PY{p}{)}
\end{Verbatim}
\end{tcolorbox}

    \begin{Verbatim}[commandchars=\\\{\}]
Prepare directory "/home/stafusa/ALF/pyALF/Notebooks/Hubbard\_Square" for Monte
Carlo run.
Resuming previous run.
Run /home/stafusa/ALF/pyALF/Notebooks/ALF/Prog/Hubbard.out
Analysing Ener\_scal
Analysing Part\_scal
Analysing Pot\_scal
Analysing Kin\_scal
Analysing Den\_eq
Analysing SpinZ\_eq
Analysing Green\_eq
Analysing SpinXY\_eq
Analysing SpinT\_eq
Analysing SpinXY\_tau
Analysing SpinZ\_tau
Analysing Den\_tau
Analysing Green\_tau
Analysing SpinT\_tau
[[-29.819654   0.135667]]

Running again reduced the error from  0.232685  to  0.135667 .
    \end{Verbatim}

    \textbf{Note}: To run a fresh simulation - instead of performing a
refinement over previous run(s) - the Monte Carlo run directory should
deleted before rerunning.

    \begin{center}\rule{0.5\linewidth}{\linethickness}\end{center}

\hypertarget{exercises}{%
\subsection{Exercises}\label{exercises}}

\begin{enumerate}
\def\labelenumi{\arabic{enumi}.}
\tightlist
\item
  Rerun once again and check the new improvement in precision.
\item
  Look at a few other observables (\texttt{sim.analysis()} outputs the
  names of those available).
\item
  Change the lattice size by adding, e.g., \texttt{"L1":\ 4,} and
  \texttt{"L2":\ 1,} to the simulation parameters definitions of
  \texttt{sim} (step 2).
\end{enumerate}


    % Add a bibliography block to the postdoc
    
    
    
